% Abhakseite „Das kann ich“ – Zone nur im Gesamt (\mitzone)
%% TODO Ich-kann-Sätze: Platzhalter wie in den Aufgabendateien
\begin{abhakseite}
\ifdefined\mitzone
\abhakgruppe{Kennst du schon}
\abhak{1}{Multiplizieren und Dividieren mit zehn, hundert und tausend als Kommaverschiebung, auch bei Dezimalzahlen}
\abhak{2}{Stellenwerttafel lesen und schreiben, Nullen an der richtigen Stelle ergänzen (null Komma null sechs Meter; eintausendfünfhundert Milliliter)}
\abhak{3}{Dezimalzahlen mit einer Zahl malnehmen und durch eine Zahl teilen, auch schriftlich oder mit dem Taschenrechner}
\abhak{4}{Dezimalzahlen vergleichen und ordnen (stellenweise, Nullen anhängen)}
\abhak{5}{Skala am Lineal und an der Uhr lesen (Untereinheiten)}
\abhak{6}{Längen im Koordinatensystem als Differenz von Koordinaten lesen (Abstand zweier Schnittstellen, Höhe zwischen Tiefpunkt und Scheitel, Funktionswert als Höhe), Flächeninhalte in Flächeneinheiten aus dem Integral oder der Figur und Volumina aus Querschnitt mal Länge}
\abhak{7}{Bruchteil einer Größe bilden (ein Viertel von einem Kilogramm; die Hälfte einer Stunde)}
\abhak{8}{Quadrat und dritte Potenz einer Länge mit Einheit (zwanzig Zentimeter zum Quadrat sind vierhundert Quadratzentimeter)}
\abhak{9}{Prozentwert einer Größe berechnen (zehn Prozent eines Volumens)}
\abhak{10}{Formel nach einer Größe umstellen (aus einer Verhältnisgleichung die gesuchte Größe isolieren)}
\abhak{11}{Fläche eines Rechtecks und eines Dreiecks berechnen (für den Materialbedarf)}
\abhak{12}{Multiplizieren und Dividieren mit zehn, hundert und tausend als Kommaverschiebung, auch bei Dezimalzahlen – Fehler finden}
\abhak{13}{Multiplizieren und Dividieren mit zehn, hundert und tausend als Kommaverschiebung, auch bei Dezimalzahlen – selbst rechnen}
\fi
\abhakgruppe{Einheit 1 · Länge, Masse, Geld}
\abhak{14}{Passt die Angabe?}
\abhak{15}{Länge, Masse und Geld umrechnen}
\abhak{16}{Maßstab im Koordinatensystem, Sek II}
\abhak{17}{Einheiten der Länge, der Masse und des Geldes der Größe nach ordnen}
\abhak{18}{Repräsentanten zuordnen (welche Angabe passt zu welchem Gegenstand)}
\abhak{19}{drei bis vier Größen ordnen}
\abhak{20}{Umrechnen mit gegebenem Faktor bei nichtmetrischen Einheiten (Fuß, Meile, Zoll; P10-Form)}
\abhak{21}{Größen mit gleicher Einheit addieren und subtrahieren}
\abhak{22}{Länge, Masse und Geld umrechnen – Fehler finden, Begründen, Anwendung, Darstellungswechsel}
\abhakgruppe{Einheit 2 · Zeit}
\abhak{23}{Zeit umrechnen und Zeitspannen berechnen}
\abhak{24}{Zeiteinheiten nennen und ordnen}
\abhak{25}{Startzeit aus Endzeit und Dauer}
\abhak{26}{Fahrplan: nächste Abfahrt, Fahrzeit}
\abhak{27}{Dauer in eine sinnvolle Einheit bringen (Sekunden in Minuten, Stunden in Tage)}
\abhak{28}{Zeit umrechnen und Zeitspannen berechnen – Fehler finden, Begründen, Anwendung, Darstellungswechsel}
\abhakgruppe{Einheit 3 · Flächen- und Volumeneinheiten}
\abhak{29}{Flächen- und Volumeneinheiten umrechnen}
\abhak{30}{Flächen- und Volumeneinheiten im Koordinatensystem, Sek II}
\abhak{31}{Flächeneinheiten ordnen (mm² bis km², a und ha einordnen)}
\abhak{32}{Repräsentanten zuordnen}
\abhak{33}{Flächen- und Volumeneinheiten umrechnen – Fehler finden, Begründen, Anwendung, Darstellungswechsel}
\abhakgruppe{Einheit 4 · Mit Größen rechnen im Sachzusammenhang}
\abhak{34}{Mit Größen rechnen im Sachzusammenhang}
\abhak{35}{Kosten und Einnahmen, Sek II}
\abhak{36}{Dichte aus Masse und Volumen (Vorrat)}
\abhak{37}{Euro und Cent in einer Rechnung}
\abhak{38}{Mit Größen rechnen im Sachzusammenhang – Fehler finden, Begründen, Anwendung}
\end{abhakseite}
